\documentclass[11pt,a4paper]{article}
\usepackage[top=1.7cm,bottom=1.5cm,left=1.5cm,right=1.5cm,headsep=8pt,footskip=16pt]{geometry}
\usepackage{enumitem}
\usepackage{amssymb}
\usepackage{array}
\usepackage[table]{xcolor}
\usepackage{fancyhdr}
\usepackage{lastpage}
\usepackage{multicol}
\usepackage{polyglossia}
\setdefaultlanguage[calendar=gregorian, numerals=maghrib]{arabic}
\setotherlanguage{english}
\IfFontExistsTF{Amiri}{%
  \newfontfamily\arabicfont[Script=Arabic, Scale=1.08]{Amiri}
}{%
  \newfontfamily\arabicfont[Script=Arabic, Scale=1.08, Path=../fonts/, Extension=.ttf,
    UprightFont=*-Regular, BoldFont=*-Bold, ItalicFont=*-Italic,
    BoldItalicFont=*-BoldItalic]{Amiri}
}
\IfFontExistsTF{AlBattar}{%
  \newfontfamily\khatt[Script=Arabic, Scale=1.15]{AlBattar}
}{%
  \newfontfamily\khatt[Script=Arabic, Scale=1.15, Path=../fonts/,
    BoldFont=ae_AlBattar.ttf]{ae_AlBattar.ttf}
}
\renewcommand{\headrulewidth}{0.8pt}
\renewcommand{\footrulewidth}{0.8pt}
\fancyhf{}
\rhead{\khatt المدرسة العليا للأساتذة ببوزريعة}
\lhead{\khatt ملحقة بومرداس}
\lfoot{\khatt الأستاذ: د. بن عبد الرزاق يوسف}
\cfoot{\thepage/\pageref{LastPage}}
\pagestyle{fancy}
\parindent 0pt
\setlength{\columnsep}{0.6cm}
\setlength{\columnseprule}{0.4pt}
\setlist[enumerate,1]{label=\textbf{\arabic*.}, leftmargin=*, itemsep=5pt, topsep=2pt}
\begin{document}

\linespread{1.08}\selectfont
\begin{center}
{\LARGE\khatt\textbf{الاختبار الثامن عشر في الذكاء الاصطناعي}}\\[4pt]
{\normalsize وضعيات إشكالية ــ أسئلة متعددة الاختيارات ــ الفصول من 1 إلى 6}
\end{center}
\vspace{-2pt}\rule{\textwidth}{0.6pt}\\[2pt]
\textbf{التعليمة:} يتضمن الاختبار خمس وضعيات إشكالية. اقرأ كل وضعية بعناية، ثم اختر التحليل أو التصرف الأنسب
من بين الاختيارات الخمسة؛ لكل وضعية إجابة صحيحة واحدة فقط.
اكتب \textbf{حرف} الإجابة الصحيحة في جدول الإجابات في آخر الاختبار.
لكل وضعية 4 نقاط (المجموع 20 نقطة).
\vspace{4pt}

\noindent
{\large\khatt\textbf{الوضعية 1}}\hfill{\small(الفصل 2)}\par\nopagebreak\vspace{3pt}
\setlength{\fboxsep}{7pt}\colorbox{gray!10}{\parbox{\dimexpr\linewidth-14pt}{بدأ أستاذ محادثة طويلة بتحديد تعليمات دقيقة: مستوى التلاميذ، وأسلوب الكتابة، وشكل التصحيح. وبعد عشرات الرسائل وإدراج نصوص طويلة، صارت الإجابات تتجاهل بعض تلك التعليمات. ويظن أن الأداة «تعبت». ويريد أن يفهم ما يحدث فعلًا ليتجنبه، لأنه يعتمد على المحادثات الطويلة في تحضير دروسه كل أسبوع. ما التفسير الأدق والإجراء الأنسب؟}}\par\nopagebreak
\begin{itemize}[leftmargin=1.0cm, labelsep=6pt, itemsep=3pt, topsep=4pt, parsep=0pt]
  \item[\textbf{أ)}] للنموذج حدّ لما يراعيه من السياق، والحل نسخ المحادثة كلها بكل نصوصها الطويلة في محادثة جديدة، لأن نقل كل شيء كما هو يعيد مراعاة التعليمات الأولى بالضرورة، ولو تجاوز طولها ما يستطيع النموذج مراعاته.
  \item[\textbf{ب)}] السبب أن النصوص المدرجة طويلة، فيحذف آخر نص أدرجه ويواصل، لأن الحذف من آخر المحادثة يعيد للنموذج تعليماته الأولى تلقائيًا، دون تذكير أو محادثة جديدة.
  \item[\textbf{ج)}] للنموذج حدّ لما يستطيع مراعاته من سياق المحادثة، ومع طولها قد تضعف مراعاته للتعليمات الأولى؛ فيذكّره بها، أو يفتح محادثة جديدة يضع التعليمات في بدايتها مع ما يلزم فقط.
  \item[\textbf{د)}] الأداة تتعب فعلًا كالإنسان بعد عمل طويل، فيتركها ساعات حتى تستريح، ثم يواصل المحادثة نفسها، فتعود إلى مراعاة كل التعليمات الأولى دون حاجة إلى تذكيرها بها، كما يحدث مع أي جهاز يعمل طويلًا.
  \item[\textbf{هـ)}] النموذج قرر تجاهل التعليمات لأنه رأى أنها غير مناسبة، فالأدوات تراجع تعليمات المستخدم وتغيّرها بحسب تقديرها، فيقبل الأستاذ الإجابات الجديدة لأنها أفضل من تعليماته.
\end{itemize}
\par\vspace{6pt}
\noindent
{\large\khatt\textbf{الوضعية 2}}\hfill{\small(الفصل 3)}\par\nopagebreak\vspace{3pt}
\setlength{\fboxsep}{7pt}\colorbox{gray!10}{\parbox{\dimexpr\linewidth-14pt}{تبحث هيئة تحرير مجلة المؤسسة عن شعار قصير بالعربية، فطلبت الأستاذة المشرفة من الأداة «اقترح شعارًا»، فجاء شعار واحد عادي. وتريد الهيئة خيارات متنوعة تناقشها، تعكس قيم المؤسسة وتكون سهلة الحفظ. وموعد طباعة العدد الأول بعد أسبوعين، وأعضاء الهيئة من التلاميذ يريدون المشاركة في الاختيار. أي طلب أنسب؟}}\par\nopagebreak
\begin{itemize}[leftmargin=1.0cm, labelsep=6pt, itemsep=3pt, topsep=4pt, parsep=0pt]
  \item[\textbf{أ)}] «اقترح عشرة شعارات لا تتجاوز خمس كلمات، تعكس قيم المعرفة والانتماء والإبداع، متنوعة الأسلوب بين السجع والاستعارة والعبارة المباشرة، مع تعليل موجز لكل منها»، ثم تناقشها الهيئة وتختار وتعدّل.
  \item[\textbf{ب)}] «اقترح مئة شعار طويل مفصّل يشرح قيم المؤسسة كلها»، لأن الشعار كلما طال كان أوضح في التعبير عن القيم، والكثرة تضمن وجود الشعار المناسب دون تحديد معايير.
  \item[\textbf{ج)}] «اقترح أفضل شعار ممكن في العالم»، لأن طلب الأفضل يجعل الأداة تقدّم أقوى اقتراح لديها، والشعار الواحد الممتاز أنفع للهيئة من خيارات كثيرة متفاوتة تضيّع وقت النقاش، وتوفر عليها عناء المقارنة.
  \item[\textbf{د)}] «اقترح عشرة شعارات»، ثم تعتمد الهيئة أول شعار في القائمة دون نقاش، لأن الأداة ترتّب اقتراحاتها حسب الجودة دائمًا، فيكون الأول أفضلها بالضرورة.
  \item[\textbf{هـ)}] «اقترح عشرة شعارات تعكس قيم المؤسسة»، ثم تعتمد الشعار الذي تفضّله الأداة عند سؤالها، لأن الأداة أقدر من الهيئة على الحكم على سهولة الحفظ ووقع الشعار في الجمهور، والهيئة تكتفي بالموافقة على اختيارها.
\end{itemize}
\par\vspace{6pt}
\noindent
{\large\khatt\textbf{الوضعية 3}}\hfill{\small(الفصل 4)}\par\nopagebreak\vspace{3pt}
\setlength{\fboxsep}{7pt}\colorbox{gray!10}{\parbox{\dimexpr\linewidth-14pt}{أخبر تلميذ أستاذته أنه يتحدث يوميًا مع روبوت محادثة عن مشكلات عائلية تقلقه، وأنه «يفهمه أكثر من الناس». وتلاحظ عليه انعزالًا متزايدًا عن زملائه. وتريد التصرف بما يحفظ مصلحته دون أن تحرجه. وهي ليست مختصة في الإرشاد النفسي، وتخشى أن يفقد التلميذ ثقته بها إن تصرفت بطريقة غير مناسبة. ما التصرف الأنسب؟}}\par\nopagebreak
\begin{itemize}[leftmargin=1.0cm, labelsep=6pt, itemsep=3pt, topsep=4pt, parsep=0pt]
  \item[\textbf{أ)}] لا تتدخل أبدًا، لأن استعمال التلميذ للأدوات خارج القسم شأن خاص لا علاقة للمدرسة به، والانعزال عن الزملاء مرحلة عابرة لا تستدعي أي توجيه، ولا يحق للأستاذة أن تسأله عن ذلك.
  \item[\textbf{ب)}] تقرأ محادثاته مع الروبوت دون علمه لتعرف مشكلاته بدقة، ثم تعرضها على زملائها، لأن معرفة التفاصيل ضرورية لمساعدته، ولو تم ذلك دون إذنه.
  \item[\textbf{ج)}] تطلب منه حذف التطبيق فورًا أمام زملائه، لأن المنع الحازم العلني يحمي التلميذ، والانعزال سيختفي تلقائيًا عندما يتوقف عن محادثة الروبوت.
  \item[\textbf{د)}] تشجعه على الاستمرار مع الروبوت وحده، لأن الأداة متاحة دائمًا ولا تحكم عليه، وهي أنسب من المختصين في معالجة المشكلات العائلية، فلا داعي لإشراك أي طرف آخر، فهي تحفظ أسراره ولا تفشيها.
  \item[\textbf{هـ)}] تستمع إليه باهتمام دون سخرية أو لوم، وتوضح أن الأداة قد تبدو متفهمة لكنها لا تعوّض الدعم الإنساني والمختصين، وتوجّهه بلطف إلى مستشار التوجيه أو الأخصائي وفق إجراءات المؤسسة.
\end{itemize}
\par\vspace{6pt}
\noindent
{\large\khatt\textbf{الوضعية 4}}\hfill{\small(الفصل 5)}\par\nopagebreak\vspace{3pt}
\setlength{\fboxsep}{7pt}\colorbox{gray!10}{\parbox{\dimexpr\linewidth-14pt}{طلبت أستاذة من أداة تلخيص قصيدة رمزية حديثة في ثلاثة أسطر، لتقديمها للتلاميذ قبل التحليل. جاء الملخص واضحًا، لكنه فسّر الرموز تفسيرًا واحدًا محددًا، وقدّمه على أنه «معنى القصيدة»، مع أن النقاد يختلفون في تأويلها. وهي تريد أن يبني التلاميذ قراءاتهم الخاصة للنص قبل الاطلاع على آراء النقاد. ما التقدير الأنسب؟}}\par\nopagebreak
\begin{itemize}[leftmargin=1.0cm, labelsep=6pt, itemsep=3pt, topsep=4pt, parsep=0pt]
  \item[\textbf{أ)}] الملخص سليم، لأن لكل قصيدة معنى واحدًا صحيحًا، والأداة حدّدته بعد تحليل آراء النقاد كلها، فتقدّمه للتلاميذ على أنه المعنى المعتمد قبل التحليل لتوفير الوقت.
  \item[\textbf{ب)}] تلخيص النص الرمزي يُسقط تعدد التأويل حين يقدّم قراءة واحدة على أنها المعنى؛ فلا تقدّم الملخص حقيقةً قبل التحليل، بل قراءةً ممكنة يناقشها التلاميذ ويقارنونها بقراءاتهم وقراءات النقاد.
  \item[\textbf{ج)}] تحذف تفسير الرموز من الملخص وتقدّم الباقي على أنه معنى القصيدة، لأن المشكلة في الرموز وحدها، أما بقية الملخص فتمثّل المعنى الوحيد الذي يتفق عليه كل النقاد، فيبدأ التلاميذ التحليل من أرضية مشتركة.
  \item[\textbf{د)}] تتجنب تقديم أي قراءة للقصيدة قبل التحليل أو بعده، لأن تعدد التأويل يعني أن القصيدة بلا معنى، فلا فائدة من مناقشة قراءات التلاميذ أو النقاد أو الأداة، ولو اختلفت آراؤهم فيها.
  \item[\textbf{هـ)}] تطلب ملخصًا أطول بعشرة أسطر وتقدّمه على أنه المعنى النهائي، لأن طول الملخص يضمن إدراج كل التأويلات الممكنة، فلا يحتاج التلاميذ إلى مناقشتها أو مقارنتها.
\end{itemize}
\par\vspace{6pt}
\noindent
{\large\khatt\textbf{الوضعية 5}}\hfill{\small(الفصل 6)}\par\nopagebreak\vspace{3pt}
\setlength{\fboxsep}{7pt}\colorbox{gray!10}{\parbox{\dimexpr\linewidth-14pt}{تريد أستاذة إعداد إملاء مسجّل بأداة تحويل النص إلى صوت، يستعمله التلاميذ في البيت للتدرب على كتابة الهمزة. جاءت القراءة سريعة متصلة، دون وقفات كافية بين الجمل، ونُطقت الهمزات بطرق متفاوتة. وتريد إرسال التسجيل إلى التلاميذ وأوليائهم في نهاية الأسبوع، واستعماله مع أقسام أخرى. ما الإعداد الأنسب لإنتاج إملاء صالح؟}}\par\nopagebreak
\begin{itemize}[leftmargin=1.0cm, labelsep=6pt, itemsep=3pt, topsep=4pt, parsep=0pt]
  \item[\textbf{أ)}] ترسل التسجيل كما هو، لأن القراءة السريعة تدرّب التلاميذ على التركيز، والتفاوت في نطق الهمزات يعكس تنوع اللهجات، فلا يؤثر في إتقان قواعد كتابة الهمزة، كما يحدث في الحياة الواقعية.
  \item[\textbf{ب)}] تقسّم النص إلى مقاطع قصيرة مع وقفات، ثم ترسل التسجيل دون الاستماع إليه، لأن التقسيم وحده يصحح نطق الهمزات ويضبط السرعة تلقائيًا في كل الأدوات الصوتية، فلا تحتاج الأستاذة إلى أي جهد إضافي.
  \item[\textbf{ج)}] تحذف الهمزات من النص قبل إدخاله، لأن الأداة تخطئ فيها، فيأتي النطق موحدًا، ويكتب التلاميذ الهمزات بأنفسهم من دون أن يسمعوها في التسجيل.
  \item[\textbf{د)}] تضبط النص بالشكل في المواضع الملتبسة، وتقسّمه إلى مقاطع قصيرة مع وقفات مقصودة، وتخفّض سرعة القراءة إن أمكن، ثم تستمع إلى التسجيل كاملًا لتتأكد من سلامة نطق الهمزات قبل إرساله.
  \item[\textbf{هـ)}] تتخلى عن الإملاء المسجّل، لأن الأدوات الصوتية لا يمكن ضبط سرعتها أو وقفاتها أو نطقها بأي طريقة، والإملاء لا يصح إلا بصوت الأستاذة مباشرة في القسم.
\end{itemize}
\par\vspace{6pt}
\begin{center}
{\khatt\textbf{جدول الإجابات}}\\[3pt]
\renewcommand{\arraystretch}{1.6}
\begin{tabular}{|>{\centering\arraybackslash}m{1.9cm}|*{5}{>{\centering\arraybackslash}m{1.2cm}|}}
\hline
\textbf{رقم السؤال} & \textbf{1} & \textbf{2} & \textbf{3} & \textbf{4} & \textbf{5} \\ \hline
\textbf{الإجابة} &  &  &  &  &  \\[8pt] \hline
\end{tabular}

\vspace{8pt}
\renewcommand{\arraystretch}{2.2}
\begin{tabular}{|>{\centering\arraybackslash}m{7.5cm}|>{\centering\arraybackslash}m{3.2cm}|>{\centering\arraybackslash}m{3.2cm}|}
\hline
\textbf{الاسم واللقب} & \textbf{الفوج} & \textbf{العلامة} \\ \hline
 & & \hfill /20 \\ \hline
\end{tabular}
\end{center}
\end{document}
